\documentclass[10pt,a4paper]{amsart}
\usepackage[margin=19mm]{geometry}
\usepackage{amsmath,amssymb,mathtools}
\usepackage[scr=rsfs,cal=euler]{mathalfa}
\usepackage{xfrac}
\usepackage[unicode,hidelinks,bookmarks=false]{hyperref}
\newcommand{\ie}{i.e.\ }
\newcommand{\cf}{cf.\ }
\newcommand{\resp}{respectively\ }
\newcommand{\Subsection}{\subsection}
\providecommand{\nobreakdash}{\nobreak\hbox{-}}
\setlength{\emergencystretch}{2em}
\multlinegap0pt
\usepackage{amssymb}
\newtheorem{theorem}{Theorem}
\title{On the special values of Berndt's generalized $L$-function}
\author{Yuan He}
\date{Source: arXiv:2608.29636v1 (CC BY 4.0)}
\begin{document}
\maketitle

\noindent\emph{Source and licence.} On the special values of Berndt's generalized $L$-function. Version arXiv:2608.29636v1 (CC BY 4.0), distributed by arXiv under the Creative Commons Attribution 4.0 International licence, http://creativecommons.org/licenses/by/4.0/, as recorded in the arXiv licence metadata for this exact version and shown on the abstract page https://arxiv.org/abs/2608.29636v1. Full licence notice: \texttt{licenses/arxiv-2608.29636-CC-BY-4.0.txt}.\par\smallskip

\noindent\textit{Editorial scope.} Counted unit: the article's theorem thm3.1 (source0.tex lines 477--483). The article's complete original proof of it is delivered with it (source0.tex lines 485--596, one \texttt{proof} environment of 112 lines). No mathematics is written, repaired or re-derived: the statement and the proof are the article's own lines, in the article's own notation, and no retained line is substituted or re-flowed.  Retained uncounted as the source-local context the proof needs: lines 80--82; lines 106--109; lines 162--164; lines 253--255; lines 373--375; lines 381--383; lines 428--430.  Nothing was omitted from any retained span, so the package carries no omission record.  No retained line contains a citation, so the wrapper carries no bibliography and no bibliography entry.  No retained line contains a figure environment, an image-inclusion command or drawing code, so no image is delivered and the package declares no asset dependency.  Editorial formatting only, in addition: an AMS \texttt{amsart} preamble that restates the article's own declarations and macros used by the retained lines, namely the environments \texttt{theorem} on separate counters, together with the standard packages those lines need.  The retained environments print in this document as Theorem 1 in order of appearance, whereas the article prints the same items with the numbering of its own full document.  The complete native source of the article as distributed in the arXiv e-print for this exact version is bundled unchanged as \texttt{sources/208872/source0.tex}.
\begin{equation}\label{eq1.11}
L(1-s,\chi)=\frac{q^{s-1}\Gamma(s)}{(2\pi)^{s}}\bigl(e^{-\frac{\pi \mathrm{i}s}{2}}+\chi(-1)e^{\frac{\pi \mathrm{i}s}{2}}\bigl)G(1,\chi)L(s,\overline{\chi}),
\end{equation}

\begin{equation}  \label{eq1.15}
\frac{te^{xt}}{\lambda e^{t}-1}=\sum_{n=0}^{\infty}
\beta_{n}(x,\lambda)\frac{t^{n}}{n!}\quad(|t+\log\lambda|<2\pi).
\end{equation}

\begin{equation}\label{eq2.5}
\beta_{n,\chi}(x,\lambda)=q^{n-1}\sum_{r=1}^{q}\chi(r)\lambda^{r}\beta_{n}\biggl(\frac{x+r}{q},\lambda^{q}\biggl)\quad(n\in\mathbb{N}_{0}).
\end{equation}

\begin{equation}\label{eq2.16}
\lambda\beta_{n}(x+1,\lambda)-\beta_{n}(x,\lambda)=nx^{n-1}\quad(n\in\mathbb{N},\lambda\in\mathbb{C}\setminus\{0\}),
\end{equation}

\begin{equation}\label{eq2.34}
\beta_{n,\chi}(\{x\}_{\chi},\lambda)=q^{n-1}\sum_{k=1}^{q}\chi(k)\lambda^{k-q[\frac{x+k}{q}]}\beta_{n}\biggl(\biggl\{\frac{x+k}{q}\biggl\},\lambda^{q}\biggl),
\end{equation}

\begin{equation}\label{eq2.35}
\overline{\beta}_{n}(x,\lambda)=\lambda^{-[x]}\beta_{n}(\{x\},\lambda)+\frac{1}{2}\delta_{1,n}\delta_{\mathbb{Z}}(x)\lambda^{-x}\quad(n\in\mathbb{N}_{0}),
\end{equation}

\begin{equation}\label{eq2.40}
\beta_{n,\chi}(\{x\}_{\chi},\lambda)=\beta_{n,\chi}(x,\lambda)-n\underset{\substack{0\leqslant k\leqslant x\\ k\in\mathbb{Z}}}{\sum}\chi(-k)\lambda^{-k}(x-k)^{n-1}.
\end{equation}

\begin{theorem}\label{thm3.1} Let $n,q\in\mathbb{N}$, $x\in\mathbb{R}$ and $\lambda\in\mathbb{C}\setminus\{0\}$. Let $\chi$ be a Dirichlet character modulo $q$. Then
\begin{equation}\label{eq3.1}
\overline{\beta}_{n,\chi}(x,\lambda)=-\frac{q^{n-1}n!}{\lambda^{x}(2\pi \mathrm{i})^{n}}\sideset{}{'}\sum_{k=-\infty}^{+\infty}\frac{G(k,\chi)e^{\frac{2\pi\mathrm{i}kx}{q}}}{(k-\frac{q\log\lambda}{2\pi\mathrm{i}})^{n}},
\end{equation}
where the dash denotes throughout that undefined terms are excluded from the
sum, and $G(k,\chi)$ is the Gauss sum defined as in \eqref{eq1.11}.
\end{theorem}

\begin{proof}
Since $\lambda^{x}\beta_{n,\chi}(\{x\}_{\chi},\lambda)$ is of bounded variation on every finite interval, by Dirichlet-Jordan test we have
\begin{equation}\label{eq3.2}
\lambda^{x}\overline{\beta}_{n,\chi}(x,\lambda)=\sum_{k=-\infty}^{+\infty}c_{n,k}(q)e^{\frac{2\pi\mathrm{i}kx}{q}},
\end{equation}
where the Fourier coefficients $c_{n,k}(q)$ are determined by
\begin{eqnarray}\label{eq3.3}
c_{n,k}(q)&=&\frac{1}{q}\int_{0}^{q}\lambda^{t}\beta_{n,\chi}(\{t\}_{\chi},\lambda)e^{-\frac{2\pi\mathrm{i}kt}{q}}\text{d}t\nonumber\\
&=&\frac{1}{q}\int_{0}^{q}\beta_{n,\chi}(\{t\}_{\chi},\lambda)e^{(\log\lambda-\frac{2\pi\mathrm{i}k}{q})t}\text{d}t.
\end{eqnarray}
Trivially, for $q=1$, we have
\begin{eqnarray}\label{eq3.4}
c_{n,k}(1)&=&\int_{0}^{1}\beta_{n}(t,\lambda)e^{(\log\lambda-2\pi\mathrm{i}k)t}\text{d}t.
\end{eqnarray}
Note that by taking the derivative with respect to $x$ on both sides of \eqref{eq1.15}, we have
\begin{equation}\label{eq3.5}
\frac{\partial}{\partial x}\beta_{n}(x,\lambda)=n\beta_{n-1}(x,\lambda)\quad(n\in\mathbb{N}).
\end{equation}
Hence, using integration by parts on \eqref{eq3.4} and in light of \eqref{eq2.16} and \eqref{eq3.5}, we know that if $\lambda\neq1$, or if $\lambda=1$ and $k\neq0$, then
\begin{eqnarray}\label{eq3.6}
c_{n,k}(1)&=&-\frac{1}{2\pi\mathrm{i}k-\log\lambda}\bigl(\lambda\beta_{n}(1,\lambda)-\beta_{n}(0,\lambda)-nc_{n-1,k}(1)\bigl)\nonumber\\
&=&-\frac{1}{2\pi\mathrm{i}k-\log\lambda}\bigl(\delta_{1,n}-nc_{n-1,k}(1)\bigl).
\end{eqnarray}
Also, from \eqref{eq1.15}, we have
\begin{equation*}
(\lambda e^{t}-1)\sum_{n=0}^{\infty}\beta_{n}(x,\lambda)\frac{t^{n}}{n!}=te^{xt}.
\end{equation*}
Applying the Taylor series of $e^{t}$ to both sides of the above identity, we find that
\begin{equation}\label{eq3.7}
\beta_{0}(x,\lambda)=\begin{cases}
1,  &\lambda=1,\\
0,  &\lambda\not=1,
\end{cases}
\end{equation}
and
\begin{equation}\label{eq3.8}
\beta_{1}(x,\lambda)=\begin{cases}
x-\frac{1}{2},  &\lambda=1,\\
\frac{1}{\lambda-1},  &\lambda\not=1.
\end{cases}
\end{equation}
It follows from \eqref{eq3.7} that if $\lambda\neq1$, or if $\lambda=1$ and $k\neq0$, then
\begin{equation*}
c_{0,k}(1)=0,
\end{equation*}
which together with \eqref{eq3.6} yields, for $\lambda\neq1$ or ($\lambda=1$ and $k\neq0$),
\begin{equation}\label{eq3.9}
c_{n,k}(1)=-\frac{n!}{(2\pi\mathrm{i}k-\log\lambda)^{n}}.
\end{equation}
Because for $\lambda=1$, by \eqref{eq2.16} and \eqref{eq3.5} we have
\begin{equation}\label{eq3.10}
c_{n,0}(1)=\int_{0}^{1}\beta_{n}(t,1)\text{d}t=0\quad(n\in\mathbb{N}),
\end{equation}
we thus see, by inserting \eqref{eq3.9} and \eqref{eq3.10} into \eqref{eq3.2}, that
\begin{equation}\label{eq3.11}
\overline{\beta}_{n}(x,\lambda)=-\frac{n!}{\lambda^{x}(2\pi \mathrm{i})^{n}}\sideset{}{'}\sum_{k=-\infty}^{+\infty}\frac{e^{2\pi\mathrm{i}kx}}{(k-\frac{\log\lambda}{2\pi\mathrm{i}})^{n}},
\end{equation}
which means Theorem \ref{thm3.1} holds for $q=1$.

We next consider the case $q\geqslant2$. Taking the derivative with respect to $x$ on both sides of \eqref{eq3.11}, we obtain that for $n\geqslant2$ and $x\in\mathbb{R}$,
\begin{equation*}
\frac{\partial}{\partial x}\overline{\beta}_{n}(x,\lambda)=n\overline{\beta}_{n-1}(x,\lambda).
\end{equation*}
From this, \eqref{eq2.34} and \eqref{eq2.35}, we see that
\begin{equation}\label{eq3.12}
\frac{\partial}{\partial x}\beta_{n,\chi}(\{x\}_{\chi},\lambda)=n\beta_{n-1,\chi}(\{x\}_{\chi},\lambda),
\end{equation}
for $n\geqslant2$ and $x\in\mathbb{R}$, except when $n=2$ and $x$ belongs to a certain set of measure zero (namely, the set of $x$ for which there exists $k\in\{1,2,\ldots,q\}$ such that $q\mid(x+k)$). Using integration by parts in \eqref{eq3.3} and taking \eqref{eq3.12} into account, we know that if $2\pi\mathrm{i}k-q\log\lambda\neq0$, then
\begin{eqnarray}\label{eq3.13}
c_{n,k}(q)=\frac{qn}{2\pi\mathrm{i}k-q\log\lambda}c_{n-1,k}(q)\quad(n\geqslant2).
\end{eqnarray}
Clearly, from \eqref{eq2.5}, \eqref{eq2.40} and \eqref{eq3.8}, it follows that if $\lambda^{q}=1$, then for $t\in\mathbb{R}$,
\begin{equation}\label{eq3.14}
\beta_{1,\chi}(\{t\}_{\chi},\lambda)=\frac{t}{q}\sum_{r=1}^{q}\chi(r)\lambda^{r}+\sum_{r=1}^{q}\chi(r)\lambda^{r}B_{1}\biggl(\frac{r}{q}\biggl)-\underset{\substack{0\leqslant r\leqslant t\\ r\in\mathbb{Z}}}{\sum}\chi(-r)\lambda^{-r},
\end{equation}
and if $\lambda^{q}\neq1$, then for $t\in\mathbb{R}$,
\begin{equation}\label{eq3.15}
\beta_{1,\chi}(\{t\}_{\chi},\lambda)=\frac{1}{\lambda^{q}-1}\sum_{r=1}^{q}\chi(r)\lambda^{r}-\underset{\substack{0\leqslant r\leqslant t\\ r\in\mathbb{Z}}}{\sum}\chi(-r)\lambda^{-r}.
\end{equation}
Inserting \eqref{eq3.14} and \eqref{eq3.15} into \eqref{eq3.3}, respectively, and then using integration by parts, we find that for $2\pi\mathrm{i}k-q\log\lambda\neq0$,
\begin{eqnarray}\label{eq3.16}
c_{1,k}&=&-\frac{1}{2\pi\mathrm{i}k-q\log\lambda}\sum_{r=1}^{q}\chi(r)\lambda^{r}\nonumber\\
&&-\frac{1}{q}\int_{0}^{q}\underset{\substack{0\leqslant r\leqslant t\\ r\in\mathbb{Z}}}{\sum}\chi(-r)\lambda^{-r}e^{(\log\lambda-\frac{2\pi\mathrm{i}k}{q})t}\text{d}t.
\end{eqnarray}
Observe that for $2\pi\mathrm{i}k-q\log\lambda\neq0$,
\begin{eqnarray}\label{eq3.17}
&&\int_{0}^{q}\underset{\substack{0\leqslant r\leqslant t\\ r\in\mathbb{Z}}}{\sum}\chi(-r)\lambda^{-r}e^{(\log\lambda-\frac{2\pi\mathrm{i}k}{q})t}\text{d}t\nonumber\\
&&\qquad=\chi(0)\lambda^{0}\int_{0}^{1}e^{(\log\lambda-\frac{2\pi\mathrm{i}k}{q})t}\text{d}t\nonumber\\
&&\qquad\quad+\bigl(\chi(0)\lambda^{0}+\chi(-1)\lambda^{-1}\bigl)\int_{1}^{2}e^{(\log\lambda-\frac{2\pi\mathrm{i}k}{q})t}\text{d}t\nonumber\\
&&\qquad\quad+\bigl(\chi(0)\lambda^{0}+\chi(-1)\lambda^{-1}+\chi(-2)\lambda^{-2}\bigl)\int_{2}^{3}e^{(\log\lambda-\frac{2\pi\mathrm{i}k}{q})t}\text{d}t\nonumber\\
&&\qquad\quad+\cdots\nonumber\\
&&\qquad\quad+\bigl(\chi(0)\lambda^{0}+\chi(-1)\lambda^{-1}+\cdots+\chi\bigl(-(q-1)\bigl)\lambda^{-(q-1)}\bigl)\int_{q-1}^{q}e^{(\log\lambda-\frac{2\pi\mathrm{i}k}{q})t}\text{d}t\nonumber\\
&&\qquad=\sum_{r=0}^{q-1}\chi(-r)\lambda^{-r}\int_{r}^{q}e^{(\log\lambda-\frac{2\pi\mathrm{i}k}{q})t}\text{d}t\nonumber\\
&&\qquad=-\frac{q}{2\pi\mathrm{i}k-q\log\lambda}\sum_{r=0}^{q-1}\chi(-r)\lambda^{-r}\bigl(\lambda^{q}-\lambda^{r}e^{-\frac{2\pi\mathrm{i}kr}{q}}\bigl)\nonumber\\
&&\qquad=-\frac{q}{2\pi\mathrm{i}k-q\log\lambda}\biggl(\sum_{r=1}^{q}\chi(r)\lambda^{r}-\sum_{r=1}^{q}\chi(r)e^{\frac{2\pi\mathrm{i}kr}{q}}\biggl).
\end{eqnarray}
So from \eqref{eq3.16} and \eqref{eq3.17}, we obtain that for $2\pi\mathrm{i}k-q\log\lambda\neq0$,
\begin{equation}\label{eq3.18}
c_{1,k}=-\frac{G(k,\chi)}{2\pi\mathrm{i}k-q\log\lambda}.
\end{equation}
It follows from \eqref{eq3.13} and \eqref{eq3.18} that if $2\pi\mathrm{i}k-q\log\lambda\neq0$, then for $n\in\mathbb{N}$,
\begin{eqnarray}\label{eq3.19}
c_{n,k}(q)=-\frac{q^{n-1}n!G(k,\chi)}{(2\pi\mathrm{i}k-q\log\lambda)^{n}}.
\end{eqnarray}
On the other hand, if $2\pi\mathrm{i}k-q\log\lambda=0$, then by using integration by parts, we know from \eqref{eq2.34}, \eqref{eq3.3} and \eqref{eq3.12} that for $n\in\mathbb{N}$,
\begin{eqnarray}\label{eq3.20}
c_{n,k}(q)&=&\frac{1}{q}\int_{0}^{q}\beta_{n,\chi}(\{t\}_{\chi},\lambda)\text{d}t\nonumber\\
&=&\frac{\beta_{n+1,\chi}(\{q\}_{\chi},\lambda)-\beta_{n+1,\chi}(\{0\}_{\chi},\lambda)}{q(n+1)}\nonumber\\
&=&0.
\end{eqnarray}
Therefore, inserting \eqref{eq3.19} and \eqref{eq3.20} into \eqref{eq3.2} yields \eqref{eq3.1}, which finishes the proof of Theorem \ref{thm3.1}.
\end{proof}

\end{document}
